\chapter[Normalization and Ordinal Analysis]{Normalization, Proof Terms, and Ordinal Analysis}
\label{ch:normalization-curry-howard}

\section{Natural-deduction normalization}
\label{sec:normalization-curry-howard-01}

The normalization theory of natural deduction developed in this chapter
originates with Prawitz \cite{Prawitz1965}.

\begin{definition}[Detour]
A detour is a segment of a natural-deduction derivation in which a formula is
introduced by an introduction rule and is immediately used as the major
premise of the matching elimination rule; the \emph{major premise} of an
elimination rule is the premise exhibiting the connective or quantifier
that the rule eliminates.  A derivation is normal when it has
no detour and satisfies the required permutation conditions for disjunction
and existential elimination.
\end{definition}

For implication, the principal reduction has the schematic form
\[
  \dfrac{
    \dfrac{\begin{matrix}{[A]}^i\\ \vdots\,D\\ B\end{matrix}}
          {A\to B}\;\to I_i
    \qquad
    \begin{matrix}\vdots\,E\\ A\end{matrix}
  }{B}\;\to E
  \quad\leadsto\quad
  \begin{matrix}\vdots\,E\\ A\\ \vdots\,D\\ B\end{matrix}\;,
\]
where \(D[E/A^i]\) substitutes the derivation \(E\) for every discharged
occurrence carrying label \(i\).  This is substitution of derivations, not term
substitution inside formulas.

For the theorem below, let \(\mathsf{NJ}_{\mathrm v}\) be the finite-tree
natural-deduction presentation of \cref{ch:prop-deduction,ch:fol-deduction}
with the rules
\(\land I,\land E_1,\land E_2,\lor I_1,\lor I_2,\lor E,\to I,\to E,
\bottom E,\forall I,\forall E,\exists I,\exists E\).  It omits
\(\mathsf{RAA}\), the primitive-negation rules, and the equality rules, and
uses \(\neg A\defeq A\to\bottom\).  Its quantifier substitutions are the total
capture-avoiding substitutions fixed earlier. Its discharge classes may be
empty: \(\to I\), either branch of \(\lor E\), and the minor premise of
\(\exists E\) may discharge zero occurrences. In the nonempty case they
discharge exactly the occurrences carrying the indicated label. All other
discharge and eigenvariable conditions remain those of the displayed rules.
This explicitly defined extension differs from the nonvacuous-discharge
presentation on Sheet~\ref{sheet:nj}; normalization here concerns
\(\mathsf{NJ}_{\mathrm v}\). Empty classes allow a reduction that erases
an unused branch to leave its surrounding discharge applications well formed.

\begin{theorem}[Normalization with empty discharge classes]
Every derivation in \(\mathsf{NJ}_{\mathrm v}\) reduces to a normal derivation with
the same conclusion and no undischarged assumption that was not already
undischarged in the original derivation.  Equivalently, after identifying
alpha-equivalent formulas, the resulting set of undischarged assumption
formulas is a subset of the original set.
\end{theorem}

\paragraph{Proof.}
Apply Corollary~\ref{thm:p07-nj-principal-permutation-normalization} to the
same formula-annotated finite derivation. Its principal and permutation
conversions are exactly the two classes used in the normality definition
above. Its complete proof gives a finite reduction with the same conclusion
and no new open assumption formula, retaining all quantifier and discharge
conditions. The continuation construction and its terminating measure are
given in Section~\ref{sec:p07-natural-deduction-normalization}, following
the separately identified source \cite{DeGroote2002NJPermutation}.

Normalization and cut elimination express related proof transformations, but
they are not one theorem.  Normalization concerns a reduction relation on
natural-deduction derivations; cut elimination concerns removal of a sequent
rule.
